\begin{abstract}
A metadata-hiding group registry --- a public-state object whose
membership is recorded as an opaque commitment and whose state
transitions are gated by zero-knowledge proofs against the prior state
--- is an emerging primitive instantiated by Tornado Cash, Semaphore,
Aztec, and the Stellar Ecosystem Proposal under recent discussion.
The literature surveys these instantiations as bundled
(curve, SNARK, chain) triples, leaving designers without a taxonomy
indexed by the deployment constraints that in practice determine which
proof system is reachable. We propose a six-axis taxonomy --- anchor
chain, SNARK system, curve and in-circuit hash, post-quantum stance,
trusted-setup model, and verification host --- under a threat model
admitting Byzantine provers, passive observers including validators,
and side-channel observers, and we analyze seven concrete configurations
across five comparison dimensions. \emph{Two} of the seven
configurations carry first-party measurements: the
Stellar~+~BLS12-381~+~PLONK production deployment (C1, five Soroban
contracts whose universal SRS is reused from the Ethereum Foundation
KZG ceremony), and a feasibility-branch
Stellar~+~hash-based-FRI-style configuration (C7) measured on the
same testnet to quantify the host-function gap empirically; the
remaining configurations cite published benchmarks date-stamped to
2026-05. Our central empirical finding is that
the FRI-style verifier on Stellar deploys today at $\sim$590$\times$
the WASM-bytecode footprint of the PLONK verifier (17.2~M vs
$\sim$29~K stroops at deploy; proof-bearing entrypoints not yet
measurable end-to-end pending a PCS-tied prover), turning the
host-function gap from a blocking constraint into an economic one.
Our central design finding is that no configuration dominates: four
deployment profiles --- cost-optimized,
post-quantum-prioritized, censorship-resistance-prioritized, and
jurisdiction-constrained --- each select a different point in the
design space, and one (the jurisdiction-constrained profile) is not
served by any surveyed configuration. We accompany the taxonomy with
a migration-path catalogue that classifies each PQ-ward transformation
as ready, blocked on chain host functions (now refined to
\emph{economically blocked}), blocked on tooling maturity, or open in
cryptographic construction; the cheap-PQ-ready-and-decentralized cell
of the design space is currently empty.
\end{abstract}
